\subsection{Extended Training and Selection Results}
\label{app:extended-training}

\paragraph{End-to-end Group DRO + Fisher-Glass frontier.}
MobileNetV3-small is fine-tuned for 10 epochs (batch 128, evaluation batch 512, seeds 17--21) with AdamW learning rate and weight decay $10^{-4}$; the final epoch is evaluated. The reported H100 configuration uses bf16 and TF32. The nuisance is training-standardized background annotation; task logits use minibatch population standard deviation plus $10^{-6}$. Groups with fewer than two observations are omitted. Ridge is $10^{-3}$, $\alpha=.95$, and projected information is floored at $10^{-8}$ before another $10^{-8}$ enters the reciprocal. Table~\ref{tab:e-frontier-performance} retains the complete frontier: stronger penalties can increase information while degrading prediction.

\begin{table}[htbp]
\centering\small
\caption{End-to-end Waterbirds Group DRO + Fisher-Glass frontier. FG entries are means $\pm$ sample standard deviations over five seeds, not confidence intervals. The Group DRO reference is the baseline from the $\lambda_{\FG}=10^{-4}$ sweep; each sweep reran its baseline, so this is not a shared paired reference for every coefficient.}
\label{tab:e-frontier-performance}\label{tab:e-frontier-information}
\begin{tabular}{@{}lrrr@{}}
\toprule
Objective / $\lambda_{\FG}$ & Avg.\ acc.\ (\%) & Worst acc.\ (\%) & Loss-CVaR$_{.95}$\\
\midrule
Group DRO reference & $85.4$ & $53.2$ & $0.925$\\
$5\times10^{-5}$ & $79.7\pm6.2$ & $56.6\pm5.8$ & $0.733\pm0.061$\\
$10^{-4}$ & $79.5\pm10.2$ & $57.7\pm8.3$ & $0.733\pm0.110$\\
$3\times10^{-4}$ & $69.9\pm9.7$ & $58.8\pm8.7$ & $0.653\pm0.064$\\
$10^{-3}$ & $62.7\pm7.9$ & $49.9\pm1.8$ & $0.710\pm0.012$\\
$3\times10^{-3}$ & $57.7\pm4.5$ & $48.6\pm6.6$ & $0.717\pm0.036$\\
$10^{-2}$ & $47.9\pm16.5$ & $22.7\pm16.7$ & $1.031\pm0.416$\\
\midrule
Objective / $\lambda_{\FG}$ & Ridge certificate & Inverse p95 & \\
\midrule
Group DRO reference & $0.259$ & $3.68$ & \\
$5\times10^{-5}$ & $0.393\pm0.059$ & $2.49\pm0.34$ & \\
$10^{-4}$ & $0.502\pm0.166$ & $2.10\pm0.58$ & \\
$3\times10^{-4}$ & $0.614\pm0.189$ & $1.71\pm0.46$ & \\
$10^{-3}$ & $0.786\pm0.124$ & $1.29\pm0.20$ & \\
$3\times10^{-3}$ & $0.817\pm0.068$ & $1.22\pm0.10$ & \\
$10^{-2}$ & $0.809\pm0.076$ & $1.23\pm0.12$ & \\
\bottomrule
\end{tabular}
\end{table}

At $\lambda_{\FG}=10^{-4}$ the certificate rises from $.259$ to $.502$ and inverse p95 falls from $3.68$ to $2.10$; at $\lambda_{\FG}=10^{-2}$ worst-group accuracy is only $22.7\%$. Historical report summaries support this table; Appendix~\ref{app:estimation-certification}'s fresh audits are a separate cohort.

\paragraph{Prospective information training.}
Use Appendix~\ref{app:prospective-intervention}'s shared model, 50-epoch schedule and seeds 22--26. Backbone fine-tuning uses AdamW with learning rate and weight decay $10^{-4}$. FG adds coefficients $\{.003,.01,.03\}$ and a background probe trained on detached features (AdamW learning rate $.001$, default weight decay $.01$), whose scores are detached in the penalty. After one penalty-free epoch, training uses $\alpha=.75$, ridge $.001$, batch standard-deviation guard $10^{-6}$, groups of at least six observations with nuisance standard deviation at least $.05$, and a $.05$ information floor before direct inversion. This floor differs from audit stabilization.

Select the final model with largest validation certificate among candidates whose overall validation accuracy is at least ERM minus $.01$. If none qualify, maximize the certificate over all FG candidates for that seed and record the fallback. Choices are $[.01,.03,.01,.003,.01]$, with fallbacks in seeds 22 and 24. All audit quantities are frozen before intervention scoring. In the prospective replication, FG raises the validation certificate from $.289$ to $.330$, but the controlled drop is $52.9$ versus $52.6$ percentage points for ERM and the originally-correct flip rate is $57.7\%$ versus $56.4\%$. The accuracy constraint fails in two seeds and background sensitivity persists.

\paragraph{Constrained hidden-slice repair.}
The primary hidden-slice cohort of Appendix~\ref{app:discovery-results} compares loss, Fisher and combined repair without opening group metadata. Each arm adds its repair penalties to ERM, includes the unchanged model, and must preserve overall, balanced and each-class validation accuracy within $.02$ of the base. Loss coefficients are $\{.1,.3,1\}$, Fisher coefficients $\{.01,.03,.1\}$, and combined candidates use their Cartesian product. The same warm start and fresh AdamW optimizer ($\mathrm{lr}=.01$, weight decay $10^{-4}$, gradient clip 10) receive at most 1,200 steps. Calibration NLL at temperature one is checked every 25 steps, with patience 12, improvement threshold $10^{-5}$, and step-zero-eligible best-checkpoint restoration. The matched ERM continuation sets both penalties to zero. Base temperature remains fixed for information training, selection and subsequent audits.

For the frozen training slice, the dimensionless Fisher penalty is
\begin{equation}
P_{\mathrm{FG}}(w)=\frac{\max\{F_{\mathrm{init}},\varepsilon\}}{F(w)+\varepsilon},
\qquad \varepsilon=10^{-8},
\label{eq:e-dimensionless-penalty}
\end{equation}
Here $F(w)$ is the frozen slice's projected predictive information and $F_{\mathrm{init}}$ its value at the unchanged model, with shrinkage $.1$, ridge $.001$, and negative numerical information clipped to zero. Within the accuracy envelope, loss/Fisher repair minimizes frozen-slice loss/maximizes information. Combined selection uses the loss--information Pareto frontier, maximizing the smaller relative improvement; ties favor mean relative improvement, then the smaller coefficient sum normalized by grid maxima, $\lambda_{\mathrm{loss}}+10\lambda_{\mathrm{Fisher}}$, then earlier index. Fresh rediscovery uses matched gate/fold seeds, so the reported population can change.

\begin{table}[htbp]
\centering\small
\caption{Constrained repair on the primary hidden-slice cohort, seeds 17--21. Accuracy, NLL and information are means $\pm$ 95\% Student-$t$ half-widths; composition is a mean percentage. Panel two audits each arm's rediscovered CLIP-defined slice. The oracle uses group annotations.}
\label{tab:e-repair-performance}\label{tab:e-repair-slices}
\begin{tabular}{@{}lrrr@{}}
\toprule
Arm & Avg.\ acc.\ (\%) & Worst acc.\ (\%) & Unchanged\\
\midrule
Matched ERM continuation & $76.3\pm0.7$ & $57.3\pm1.3$ & ---\\
Loss repair & $76.3\pm0.6$ & $56.9\pm1.3$ & $3/5$\\
Fisher repair & $76.3\pm0.7$ & $57.3\pm1.3$ & $5/5$\\
Combined repair & $76.4\pm0.9$ & $56.7\pm1.6$ & $4/5$\\
Oracle Group DRO & $80.0\pm0.3$ & $64.8\pm0.9$ & ---\\
\midrule
Arm & Slice NLL & Projected info. & LL mass (\%)\\
\midrule
Matched ERM continuation & $0.031\pm0.017$ & $0.005\pm0.001$ & $97.0$\\
Loss repair & $0.023\pm0.017$ & $0.005\pm0.002$ & $97.8$\\
Fisher repair & $0.031\pm0.017$ & $0.005\pm0.001$ & $97.0$\\
Combined repair & $0.031\pm0.017$ & $0.005\pm0.001$ & $97.2$\\
Oracle Group DRO & $0.033\pm0.018$ & $0.005\pm0.002$ & $95.9$\\
\bottomrule
\end{tabular}
\end{table}

Fisher repair chooses the unchanged model in every seed. Paired worst-group changes versus matched continuation are $-.4$, $.0$ and $-.6$ points for loss, Fisher and combined repair, with intervals containing zero. The oracle reaches $64.8\%$ worst-group accuracy while rediscovered information remains $.005$. LL denotes landbird/land.

\paragraph{Post hoc group-free coefficient selection.}
This analysis was designed after earlier group/test outcomes were inspected. Its grid $\{0,.01,.03,.1,.3\}$ uses~\eqref{eq:e-dimensionless-penalty}, not the known-group objective's scale. Training is partitioned $70/15/15\%$ for fitting/calibration/gate discovery; a frozen train-only slice supplies the penalty. All candidates share the warm start, training seed, optimizer, 1,200-step budget, checkpoint rule and base temperature above. Official validation enters only after heads are fixed; group metadata stays sealed through blind test evaluation.

Let $\widehat R_\lambda$ be mean fixed-temperature validation NLL, $s_\lambda=\operatorname{sd}_i(\ell_{\lambda i})/\sqrt{1199}$ its within-run standard error, and $\lambda_R=\arg\min_\lambda(\widehat R_\lambda,\lambda)$ lexicographically. Here $\ell_{\lambda i}$ is the fixed-temperature validation NLL of example $i$, and $\operatorname{sd}_i$ is the sample standard deviation with divisor $1198$. Five task-stratified validation folds rediscover candidate-specific slices; $V_{k,\lambda}$ is held-out inverse predictive information and $s_{V,\lambda}$ its sample standard deviation. Selection uses
\begin{equation}
\begin{aligned}
\Lambda_{\mathrm{risk}}&=\{\lambda:\widehat R_\lambda\le\widehat R_{\lambda_R}+s_{\lambda_R}\},&
U_\lambda^{\mathrm{adj}}&=\overline V_\lambda+t_{4,.99}s_{V,\lambda}/\sqrt5,\\
\widehat\lambda_{\mathrm{hidden}}
&=\min\arg\max_{\lambda\in\Lambda_{\mathrm{risk}}}(U_\lambda^{\mathrm{adj}})^{-1},&
t_{4,.99}&=3.74695 .
\end{aligned}
\end{equation}
The risk allowance uses the minimizer's SE for all candidates. Bonferroni adjustment covers five coefficients, but overlapping gate-training folds make this a ranking score without independent-fold coverage. Accuracy screens and paired inverse-information tests are diagnostics, not selection constraints. The group-supervised comparator maximizes validation worst-group accuracy over the same $\Lambda_{\mathrm{risk}}$, breaking ties toward smaller $\lambda$.

\begin{table}[htbp]
\centering\small
\caption{Post hoc coefficient selection, five paired seeds. Accuracies are test results; information is projected predictive information in test slices defined by validation-rediscovered, frozen gates. Brackets are 95\% Student-$t$ intervals. Policies selecting identical models across these seeds share rows; the text explains both coincidences. Coefficient counts are $\lambda:\text{number of seeds}$.}
\label{tab:e-groupfree-performance}\label{tab:e-groupfree-information}
\begin{tabular}{@{}lrr@{}}
\toprule
Policy & Average accuracy (\%) & Worst accuracy (\%)\\
\midrule
Frozen ERM / matched zero & $75.5\ [74.3,76.6]$ & $55.7\ [53.3,58.2]$\\
Fixed $\lambda=.03$ & $76.6\ [76.1,77.2]$ & $51.5\ [47.1,55.8]$\\
Risk only / group-free & $77.5\ [76.5,78.6]$ & $43.1\ [33.1,53.1]$\\
Group-supervised comparator & $77.7\ [76.6,78.8]$ & $44.2\ [34.2,54.3]$\\
\midrule
Policy & Coefficient counts & Slice information\\
\midrule
Frozen ERM / matched zero & --- / $0:5$ & $0.007\ [0.005,0.009]$\\
Fixed $\lambda=.03$ & $.03:5$ & $0.007\ [0.004,0.010]$\\
Risk only / group-free & $0:1,\ .3:4$ & $0.012\ [0.006,0.019]$\\
Group-supervised comparator & $0:1,\ .1:1,\ .3:3$ & $0.011\ [0.005,0.018]$\\
\bottomrule
\end{tabular}
\end{table}

Matched zero continuation restores step zero in all five seeds, explaining equality with frozen ERM despite intervening optimization. Group-free selection equals plain NLL minimization in every seed; the information criterion does not change any choice. Versus matched continuation, paired mean changes are $+2.1$ average-accuracy points ($[+.2,+4.0]$), $+.006$ information ($[+.001,+.010]$), and $-12.6$ worst-group points ($[-21.6,-3.7]$). The group-supervised comparator opens validation annotations only after blind decisions and records are fixed and is not an upper bound. The post hoc coefficient-selection results and primary repair cohort have separate artifact identifiers (see the \hyperref[sec:artifact-availability]{artifact availability statement}), alongside extended settings and full tables.

\paragraph{Operational choice of the penalty.}
For group-free selection, we propose prespecifying an information requirement and an ordinary-risk budget, then choosing the smallest coefficient satisfying both. If no candidate qualifies, report the requirement as unmet. With neither a specified information requirement nor a group-labeled audit set, retain $\lambda_{\FG}=0$. A worst-group-accuracy objective requires group labels or a justified group proxy. This threshold-based policy requires prospective validation; the reported experiment maximizes information within an NLL eligibility set.
